\subsection{Backend}

Các công nghệ cho phần Backend của hệ thống thương mại điện tử cho bán lẻ điện thoại tích hợp trợ lý ảo AI Agent của nhóm được lựa chọn một cách cẩn trọng nhằm đáp ứng trọn vẹn định hướng kiến trúc hiện đại và tối ưu hóa hiệu suất.  

Các công nghệ như Spring Boot 3.x, FastAPI, Eclipse Vert.x và Apache Kafka được chọn lọc để đảm bảo tính đa dạng và phù hợp trong việc hiện thực hóa một hệ thống Microservice phân tán và Event-Driven Architecture. Sự kết hợp của các framework hiệu suất cao và nền tảng truyền tải sự kiện này là cơ sở để đạt được khả năng mở rộng (scalability) và linh hoạt cao.

Tiếp đó, nhóm đã quyết định tận dụng mô hình đa ngôn ngữ bằng cách kết hợp Java 21 và Python 3.10+. Java được chọn cùng với Spring Boot cho các dịch vụ cốt lõi nhờ sự mạnh mẽ và ổn định trong việc xử lý nghiệp vụ backend và giao dịch, trong khi đó, Python được dành riêng cho dịch vụ AI Agent, tận dụng tối đa hệ sinh thái chuyên sâu về trí tuệ nhân tạo và học máy. Bên cạnh đó, các công nghệ hỗ trợ quan trọng khác như Eureka (Service Discovery) và các framework Reactive (aiokafka, Vert.x) đã được tích hợp để đảm bảo tính nhất quán của dữ liệu và hiệu suất xử lý I/O vượt trội trong môi trường phân tán.

\subsubsection{Ngôn ngữ}
Hệ thống của nhóm áp dụng chiến lược đa ngôn ngữ (Polyglot Programming) nhằm tối ưu hóa hiệu suất cho từng vai trò dịch vụ. Nhóm đã chọn sử dụng Java 21 cho các Microservice cốt lõi và Python 3.10+ cho Microservice AI Agent.

\begin{itemize}
    \item \textbf{Java 21:} Là lựa chọn hàng đầu cho các dịch vụ cốt lõi, chẳng hạn như các nghiệp vụ quan trọng như quản lý người dùng và đơn hàng. Java, với ưu điểm về tính ổn định, khả năng mở rộng đã được chứng minh và hiệu suất cao trên máy ảo JVM, đảm bảo độ tin cậy cần thiết cho các tác vụ xử lý nghiệp vụ yêu cầu độ bền bỉ cao. 
    \item \textbf{Python 3.10+: } Được dành riêng để phát triển dịch vụ AI Agent. Python là ngôn ngữ tiêu chuẩn trong lĩnh vực trí tuệ nhân tạo và Học máy, chủ yếu nhờ vào hệ sinh thái thư viện phong phú (như TensorFlow, PyTorch). Do đó, việc sử dụng Python giúp nhóm tăng tốc độ phát triển, triển khai và bảo trì các mô hình AI phức tạp một cách hiệu quả hơn.
\end{itemize}

Sự kết hợp chiến lược này cho phép hệ thống tận dụng sức mạnh xử lý ổn định của Java cho các nghiệp vụ lõi, đồng thời khai thác sự linh hoạt và thư viện chuyên biệt của Python cho các tác vụ tính toán thông minh, tạo nên một giải pháp cân bằng và hiệu quả toàn diện.

\subsubsection{Các framework cơ bản và reactive (Core \& Reactive Frameworks)}

Để tối ưu hóa quy trình phát triển và hiệu năng thực thi của các dịch vụ, hệ thống đã lựa chọn các framework chính có tính năng bổ trợ cao cho Java và Python.

Các framework phát triển cốt lõi được chọn để làm nền tảng xây dựng API và logic nghiệp vụ, tập trung vào hiệu suất và tốc độ phát triển:

\begin{itemize}
    \item \textbf{Spring Boot 3.x\: } Đây là framework lý tưởng để xây dựng các RESTful API và xử lý logic nghiệp vụ nhờ khả năng tự động cấu hình (autoconfiguration) mạnh mẽ. Điều này giúp giảm thiểu đáng kể thời gian thiết lập ban đầu, cho phép nhóm tập trung hoàn toàn vào việc phát triển tính năng. Bên cạnh đó Spring Boot còn có khả năng tích hợp tốt với Docker, Kubernetes và CI/CD giúp tối ưu quy trình triển khai. Nhờ đó, Spring Boot trở thành nền tảng linh hoạt, dễ mở rộng và phù hợp cho các ứng dụng đòi hỏi hiệu năng ổn định và phát triển nhanh.
    \item \textbf{FastAPI: }: FastAPI nổi bật ở khả năng xây dựng API hiệu suất cao và hỗ trợ bất đồng bộ (Async/Await) vượt trội, tận dụng tối đa sức mạnh của Python cho các tác vụ tính toán chuyên sâu trong dịch vụ AI. Nhờ thiết kế nhẹ, cú pháp rõ ràng và tốc độ xử lý cao, FastAPI đặc biệt phù hợp cho các hệ thống yêu cầu phản hồi nhanh, xử lý AI theo thời gian thực và các dịch vụ vi mô Python hiện đại.
\end{itemize}

Bên cạnh các framework cốt lõi, việc tích hợp các khung xử lý phản ứng là thiết yếu để đảm bảo hiệu suất I/O vượt trội và duy trì tính non-blocking trong kiến trúc Event-Driven. Các công cụ này giúp các dịch vụ xử lý đồng thời một lượng lớn dữ liệu sự kiện một cách hiệu quả.

\begin{itemize}
    \item \textbf{Vert.x\: }Được sử dụng trong các dịch vụ Java, Vert.x là một toolkit bất đồng bộ và hướng sự kiện. Bằng cách này, Vert.x giúp ứng dụng xử lý một lượng lớn kết nối đồng thời với độ trễ thấp, tối ưu hóa việc sử dụng tài nguyên. Một điểm nổi bật của Vert.x là event bus nội bộ, cho phép các thành phần trong hệ thống trao đổi dữ liệu, gửi và nhận thông điệp một cách non-blocking mà không cần phụ thuộc trực tiếp vào nhau\cite{vertx_core}, từ đó giảm thiểu độ kết dính (coupling) giữa các module. Với kiến trúc nhẹ, đa ngôn ngữ và hiệu năng cao, Vert.x đặc biệt phù hợp cho các dịch vụ vi mô\cite{vertx_microservices}, hệ thống IoT và các ứng dụng cần throughput lớn cùng khả năng mở rộng linh hoạt.
    \item \textbf{Aiokafka: }Thư viện này được sử dụng trong Python. Nó tận dụng khả năng bất đồng bộ gốc của Python, đảm bảo rằng tiến trình giao tiếp sự kiện với Apache Kafka diễn ra non-blocking, duy trì hiệu năng cao cho việc truyền tải sự kiện xuyên suốt hệ thống\cite{aiokafka_perf}. Nhờ hỗ trợ cơ chế async/await, thư viện cho phép xử lý nhiều tác vụ đồng thời như publish, subscribe và consumer polling mà không làm nghẽn luồng\cite{aiokafka_perf}. Điều này đặc biệt hữu ích trong các hệ thống event-driven hoặc microservices cần throughput lớn, độ trễ thấp và khả năng mở rộng linh hoạt.
\end{itemize}

\subsubsection{Truy cập cơ sở dữ liệu và ánh xạ đối tượng – quan hệ (Database Access \& ORM)}

Để quản lý việc tương tác với cơ sở dữ liệu một cách hiệu quả và đảm bảo tính nhất quán của dữ liệu (data integrity), hệ thống áp dụng các công cụ ORM phù hợp với từng ngôn ngữ.
\begin{itemize}
    \item \textbf{Spring Data JPA (Sử dụng Hibernate):} Spring Data JPA được xác định là tiêu chuẩn trong việc truy cập dữ liệu cho các Microservice Java. Vai trò cốt lõi của framework này là giúp nhóm giảm thiểu đáng kể mã lặp lại (boilerplate code) bằng cách tự động hóa việc tạo ra các phương thức truy vấn dựa trên tên của chúng, nhờ đó làm tăng tốc độ phát triển và duy trì tính rõ ràng của mã. Bên cạnh đó, bằng cách sử dụng Hibernate làm nhà cung cấp JPA, các dịch vụ Java được hưởng lợi từ khả năng ánh xạ đối tượng-quan hệ (ORM) mạnh mẽ. Khả năng này giúp chuyển đổi dữ liệu phức tạp giữa đối tượng Java và bảng cơ sở dữ liệu một cách mượt mà và an toàn, đảm bảo tính bền vững của dữ liệu.
    \item \textbf{Milvus (Python SDK: pymilvus):} Milvus được lựa chọn làm nền tảng Vector Database chuyên dụng cho dịch vụ AI Agent, thay thế cho các ORM truyền thống. Vai trò cốt lõi của Milvus là lưu trữ và truy vấn các vector embeddings đa chiều. Việc giao tiếp với Milvus được thực hiện thông qua thư viện pymilvus (Python SDK), một công cụ chuyên biệt cho các thao tác vector. Ưu điểm chính của Milvus là khả năng thực hiện truy vấn vector hiệu suất cao, cần thiết cho các tác vụ tìm kiếm ngữ nghĩa và gợi ý sản phẩm dựa trên độ tương đồng. Điều này đảm bảo rằng các phản hồi của trợ lý ảo AI Agent diễn ra nhanh chóng và chính xác, vượt trội so với cơ sở dữ liệu quan hệ trong việc xử lý dữ liệu đa chiều.
\end{itemize}

\subsubsection{Nền tảng truyền tải sự kiện (Event Messaging Platform)}

Việc duy trì giao tiếp bất đồng bộ giữa các dịch vụ phân tán là yếu tố then chốt để đảm bảo tính sẵn sàng và khả năng mở rộng của hệ thống. Khi các dịch vụ có thể trao đổi dữ liệu mà không cần chờ đợi lẫn nhau, hệ thống sẽ giảm thiểu độ trễ, hạn chế nghẽn cổ chai và dễ dàng thích ứng với lưu lượng tăng cao. Xuất phát từ yêu cầu đó, nhóm đã lựa chọn một nền tảng chuyên biệt để xử lý luồng sự kiện giữa các thành phần trong hệ thống — một giải pháp tin cậy, ổn định và được ứng dụng rộng rãi trong kiến trúc hiện đại.

\textbf{Apache Kafka} được lựa chọn làm trung tâm dữ liệu bất biến và nền tảng truyền tải sự kiện cốt lõi, phục vụ cho toàn bộ giao tiếp không đồng bộ giữa các microservice\cite{kafka_event_streaming}. Nhờ khả năng xử lý lượng dữ liệu lớn với hiệu năng cao và tính bền vững nhờ cơ chế ghi log, Kafka đóng vai trò là kênh trao đổi thông điệp chính trong hệ thống. Trước hết, Kafka đảm bảo sự tách biệt lỏng lẻo giữa bên gửi và bên nhận, giúp mở rộng hệ thống linh hoạt bằng cách bổ sung dịch vụ mới mà không cần can thiệp vào các thành phần sẵn có. Đồng thời, Kafka không chỉ hoạt động như một message broker mà còn như một event store bền vững, lưu giữ toàn bộ chuỗi sự kiện theo thời gian. Đây là nền tảng quan trọng để hiện thực hóa mô hình event sourcing, khi các dịch vụ có thể khôi phục trạng thái bằng cách đọc lại lịch sử sự kiện sau sự cố hoặc khi cần kiểm tra dữ liệu\cite{kafka_event_sourcing}. Cuối cùng, việc triển khai Kafka theo chế độ KRaft giúp loại bỏ sự phụ thuộc vào Zookeeper, đơn giản hóa kiến trúc, giảm chi phí vận hành và góp phần nâng cao tính ổn định của cụm Kafka.

\subsubsection{Cơ chế khám phá dịch vụ (Service Discovery)}

Trong kiến trúc microservice, nơi các dịch vụ liên tục được mở rộng, triển khai và thay đổi địa chỉ mạng, một cơ chế khám phá dịch vụ là yếu tố bắt buộc để các thành phần có thể tự động tìm thấy và giao tiếp với nhau. Trong nền tảng này, Eureka được sử dụng như máy chủ đăng ký dịch vụ cho các microservice Java cốt lõi. Đây là giải pháp service discovery của Netflix, được thiết kế dành riêng cho môi trường cloud.

\textbf{Eureka} mang lại nhiều ưu điểm nổi bật. Trước hết, nó hỗ trợ cơ chế tự động đăng ký và phát hiện dịch vụ: khi một dịch vụ khởi động, nó sẽ tự động đăng ký địa chỉ mạng của mình và được tự động gỡ bỏ khi ngừng hoạt động, giúp loại bỏ việc cấu hình thủ công và đơn giản hóa đáng kể quá trình triển khai\cite{eureka_autoconfig}. Bên cạnh đó, Eureka có khả năng phục hồi tốt trước các lỗi mạng tạm thời. Ngay cả khi máy chủ Eureka không khả dụng, các dịch vụ vẫn có thể sử dụng dữ liệu lưu trong bộ nhớ đệm để truy cập thông tin về các dịch vụ khác, đảm bảo hệ thống tiếp tục vận hành mà không bị gián đoạn\cite{eureka_resilience}. 

\subsubsection{Tóm tắt công nghệ}
\begin{table}[H]
\centering
\caption{Tổng hợp các công nghệ sử dụng tại Backend}
\label{tab:backend_tech}
\begin{tabular}{|l|p{4.5cm}|p{5.5cm}|} 
\hline
\textbf{Thành phần} & \textbf{Công nghệ} & \textbf{Ghi chú/Vai trò} \\ 
\hline
Ngôn ngữ & Java 21, Python 3.10+ & Java cho các microservice chính; Python cho microservice AI. \\ 
\hline
Framework chính & Spring Boot 3.x, FastAPI & Viết REST API và xử lý business logic. \\ 
\hline
Khung Xử lý Phản ứng & Eclipse Vert.x, aiokafka & Xử lý I/O bất đồng bộ và giao tiếp Kafka hiệu suất cao. \\ 
\hline
Event Messaging & Apache Kafka (Kraft) & Giao tiếp bất đồng bộ giữa các service. \\ 
\hline
Service Discovery & Eureka & Tự động phát hiện và định danh service trong kiến trúc phân tán. \\ 
\hline
Database Access \& ORM & Spring Data JPA (Hibernate), SQLAlchemy (Python) & Ánh xạ dữ liệu giữa đối tượng (object) và cơ sở dữ liệu (DB). \\ 
\hline
\end{tabular}
\end{table}




